% !TEX root = ../prise-en-main-medical-vrp.tex
\section{Organiser votre activité}
\label{sec:organisation}

\subsection{Ordre recommandé}

Configurez vos données dans cet ordre — une seule fois, puis des ajustements ponctuels~:

\begin{enumerate}[label=\textbf{\arabic*.}, leftmargin=*]
  \item \textbf{Prestations} — types d'actes ou protocoles (\emph{consultation}, \emph{soins à domicile}, \emph{bilan}, etc.).
  \item \textbf{Structure} — lieu ou organisme de facturation (cabinet, clinique, domicile patient, centre partenaire).
  \item \textbf{Séances} — actes rattachés à une \structure{} et une \prestation{} (tarif, durée type).
  \item \textbf{Patients} — un \patient{} par parcours de soin / bénéficiaire, lié aux \seance{}s concernées.
\end{enumerate}

\subsection{Créer une prestation}

\begin{workflowstep}[Prestations]
  Menu \menu{Référentiel} → onglet \menu{Prestations} (\texttt{/program}). \menu{Ajouter une prestation}. Nommez l'acte de façon claire pour vos rapports.
\end{workflowstep}

\screenshot[keepaspectratio]{03-prestations.png}{Catalogue des prestations}

\subsection{Créer une structure}

\begin{workflowstep}[Structure]
  Onglet \menu{Structures} ou bouton depuis l'accueil. Renseignez nom, un \textbf{code} optionnel (unique dans l'entreprise), adresse de facturation, identifiants légaux si facturation directe. Un code déjà utilisé par une autre structure du même compte est refusé.
\end{workflowstep}

Chaque \structure{} concentre~: \seance{}s, documents, factures et la \textbf{préparation de facturation} mensuelle.

\subsection{Créer une séance (catalogue)}

\begin{workflowstep}[Séance]
  Depuis la fiche \structure{} → \menu{Ajouter une séance}. Liez la \prestation{}, définissez libellé et volume. Le tarif horaire se saisit en \textbf{TTC} par défaut (radio HT disponible)~; la base enregistre le HT avec 4 décimales.
\end{workflowstep}

\subsection{Créer un patient}

\begin{workflowstep}[Patient]
  Onglet \menu{Patients} (\texttt{/group}). \menu{Nouveau patient}. Un \patient{} distinct par parcours évite les confusions en facturation.
\end{workflowstep}

\screenshot[keepaspectratio]{04-patients.png}{Liste des patients}

\begin{astuce}
  L'import calendrier (onglet \menu{Calendrier} de l'agenda) mappe les libellés ICS vers une \seance{} \emph{et} un \patient{} de la \structure{} choisie. Les deux sont requis pour créer une session. Le réimport n'est fiable que si le mapping mémorisé est un patient (le morph stocké privilégie sinon la séance et ignore les événements).
\end{astuce}

\subsection{Exemple minimal}

\begin{itemize}[leftmargin=*]
  \item \structure{} «~Cabinet Dr Martin~»
  \item \prestation{} «~Consultation~» → \seance{} «~Consultation standard~» → \patient{} «~Dupont~»
  \item \prestation{} «~Soins infirmiers~» → \seance{} «~Pansement~» → \patient{} «~Bernard~»
\end{itemize}
